% GENERATED FILE. Edit canonical lessons, then run npm run generate:books.
% canonical-source-hash: fnv1a64:6cc3a5d26420544f
% canonical-lessons: MR-C550-yogayog, MR-C550-jugar, MR-C550-bhagya, MR-C550-nanephek, MR-C550-paij-lavne, MR-R550-review-544-550

\chapter{Luck and Chance}
\label{ch:mr-luck-and-chance}
\hlchaptermodality{550}

\begin{chapteropening}
\textbf{By the end of this chapter:} \emph{I can talk about luck, coincidences and bets.}

Complete the last lesson of chapter 550: I can talk about luck, coincidences and bets.
\end{chapteropening}

\section[\texorpdfstring{\mr{योगायोग} (yogāyog)}{योगायोग (yogāyog)}]{\mr{योगायोग} (yogāyog) — a coincidence}

\label{lesson:MR-C550-yogayog}



\begin{quote}
Before the new one: say the Marathi for protection, then the Marathi for reliance.
\end{quote}

\subsection*{You'll want to know: \mr{योगायोग}}
\textbf{\mr{योगायोग}} — \emph{yogāyog} — \textquotedblleft{}a coincidence\textquotedblright{}.

\begin{quote}
\mr{बाजारात} \mr{जुना} \mr{मित्र} \mr{भेटला}, \mr{हा} \mr{मोठा} \mr{योगायोग} \mr{होता}.
\end{quote}

\begin{quote}
\emph{bājārāt junā mitra bheṭlā, hā moṭhā yogāyog hotā.}
\end{quote}

\begin{quote}
Meeting an old friend in the market was a great coincidence.
\end{quote}

Masculine.

\begin{cousinweb}
\textbf{\mr{योगायोग}} (\emph{yogāyog}) is a tatsama, Sanskrit \emph{yogāyoga}, \emph{yoga}, \textquotedblleft{}joining\textquotedblright{}, with \emph{ayoga}, its opposite, glossed as fitness and unfitness (Monier-Williams). Marathi uses it for things that happen to meet.
\end{cousinweb}

\begin{grammarlens}[title={What you've built}]
One more word for saying what happens if.
\end{grammarlens}

\subsection*{Your turn}
\begin{itemize}
\raggedright
  \item \emph{Say it:} \emph{yogāyog}
  \item \emph{Say it:} \emph{yogāyog}, once more
  \item \emph{Say it:} \emph{bājārāt junā mitra bheṭlā, hā moṭhā yogāyog hotā.} — Meeting an old friend in the market was a great coincidence.
  \item \emph{Recall:} say the Marathi for protection, then the Marathi for reliance, then say \emph{yogāyog} again
\end{itemize}

\begin{tcolorbox}[breakable,colback=teal!4,colframe=teal!35!black,title={Before you move on}]
What does \mr{योगायोग} mean? (\textquotedblleft{}A coincidence\textquotedblright{}.) Say it once more.
\end{tcolorbox}
\section[\texorpdfstring{\mr{जुगार} (jugār)}{जुगार (jugār)}]{\mr{जुगार} (jugār) — gambling}

\label{lesson:MR-C550-jugar}



\begin{quote}
Before the new one: say the Marathi for reliance, then the Marathi for a coincidence.
\end{quote}

\subsection*{You'll want to know: \mr{जुगार}}
\textbf{\mr{जुगार}} — \emph{jugār} — \textquotedblleft{}gambling\textquotedblright{}.

\begin{quote}
\mr{जुगारात} \mr{त्याने} \mr{सगळे} \mr{पैसे} \mr{घालवले}.
\end{quote}

\begin{quote}
\emph{jugārāt tyāne sagḷe paise ghālavle.}
\end{quote}

\begin{quote}
He lost all his money gambling.
\end{quote}

Masculine.

\mr{जुगार} \mr{खेळणे}, to gamble.

\begin{cousinweb}
\textbf{\mr{जुगार}} (\emph{jugār}) is usually linked with Sanskrit \emph{dyūtakāra}, \textquotedblleft{}gambler\textquotedblright{}, under which CDIAL lists Prakrit \emph{juāra} and Hindi \emph{juwārī} (CDIAL 6609). CDIAL gives no Marathi form, so the path is not settled.
\end{cousinweb}

\begin{grammarlens}[title={What you've built}]
One more word for saying what happens if.
\end{grammarlens}

\subsection*{Your turn}
\begin{itemize}
\raggedright
  \item \emph{Say it:} \emph{jugār}
  \item \emph{Say it:} \emph{jugār}, once more
  \item \emph{Say it:} \emph{jugārāt tyāne sagḷe paise ghālavle.} — He lost all his money gambling.
  \item \emph{Recall:} say the Marathi for reliance, then the Marathi for a coincidence, then say \emph{jugār} again
\end{itemize}

\begin{tcolorbox}[breakable,colback=teal!4,colframe=teal!35!black,title={Before you move on}]
What does \mr{जुगार} mean? (\textquotedblleft{}Gambling\textquotedblright{}.) Say it once more.
\end{tcolorbox}
\section[\texorpdfstring{\mr{भाग्य} (bhāgya)}{भाग्य (bhāgya)}]{\mr{भाग्य} (bhāgya) — good fortune}

\label{lesson:MR-C550-bhagya}



\begin{quote}
Before the new one: say the Marathi for a coincidence, then the Marathi for gambling.
\end{quote}

\subsection*{You'll want to know: \mr{भाग्य}}
\textbf{\mr{भाग्य}} — \emph{bhāgya} — \textquotedblleft{}good fortune\textquotedblright{}.

\begin{quote}
\mr{अशी} \mr{आई} \mr{मिळणं} \mr{हे} \mr{माझं} \mr{भाग्य} \mr{आहे}.
\end{quote}

\begin{quote}
\emph{aśī āī miḷṇaṁ he mājhaṁ bhāgya āhe.}
\end{quote}

\begin{quote}
Having a mother like this is my good fortune.
\end{quote}

Neuter.

\begin{cousinweb}
\textbf{\mr{भाग्य}} (\emph{bhāgya}) is a tatsama, taken from Sanskrit \emph{bhāgya}, \textquotedblleft{}one's lot, fortune\textquotedblright{} (Monier-Williams).
\end{cousinweb}

\begin{grammarlens}[title={What you've built}]
One more word for saying what happens if.
\end{grammarlens}

\subsection*{Your turn}
\begin{itemize}
\raggedright
  \item \emph{Say it:} \emph{bhāgya}
  \item \emph{Say it:} \emph{bhāgya}, once more
  \item \emph{Say it:} \emph{aśī āī miḷṇaṁ he mājhaṁ bhāgya āhe.} — Having a mother like this is my good fortune.
  \item \emph{Recall:} say the Marathi for a coincidence, then the Marathi for gambling, then say \emph{bhāgya} again
\end{itemize}

\begin{tcolorbox}[breakable,colback=teal!4,colframe=teal!35!black,title={Before you move on}]
What does \mr{भाग्य} mean? (\textquotedblleft{}Good fortune\textquotedblright{}.) Say it once more.
\end{tcolorbox}
\section[\texorpdfstring{\mr{नाणेफेक} (nāṇephek)}{नाणेफेक (nāṇephek)}]{\mr{नाणेफेक} (nāṇephek) — a coin toss}

\label{lesson:MR-C550-nanephek}



\begin{quote}
Before the new one: say the Marathi for gambling, then the Marathi for good fortune.
\end{quote}

\subsection*{You'll want to know: \mr{नाणेफेक}}
\textbf{\mr{नाणेफेक}} — \emph{nāṇephek} — \textquotedblleft{}a coin toss\textquotedblright{}.

\begin{quote}
\mr{सामन्याआधी} \mr{कर्णधारांनी} \mr{नाणेफेक} \mr{केली}.
\end{quote}

\begin{quote}
\emph{sāmnyāādhī karṇadhārānnī nāṇephek kelī.}
\end{quote}

\begin{quote}
Before the match the captains tossed a coin.
\end{quote}

Feminine.

\mr{नाणेफेक} \mr{करणे}, to toss a coin.

\begin{cousinweb}
\textbf{\mr{नाणेफेक}} (\emph{nāṇephek}) is a Marathi compound built from \mr{नाणे}, \textquotedblleft{}a coin\textquotedblright{}, inherited from Sanskrit \emph{nāṇaka}, \textquotedblleft{}coin\textquotedblright{} (CDIAL 7050), and \mr{फेक}, \textquotedblleft{}a throw\textquotedblright{}, from \mr{फेकणे}.
\end{cousinweb}

\begin{grammarlens}[title={What you've built}]
One more word for saying what happens if.
\end{grammarlens}

\subsection*{Your turn}
\begin{itemize}
\raggedright
  \item \emph{Say it:} \emph{nāṇephek}
  \item \emph{Say it:} \emph{nāṇephek}, once more
  \item \emph{Say it:} \emph{sāmnyāādhī karṇadhārānnī nāṇephek kelī.} — Before the match the captains tossed a coin.
  \item \emph{Recall:} say the Marathi for gambling, then the Marathi for good fortune, then say \emph{nāṇephek} again
\end{itemize}

\begin{tcolorbox}[breakable,colback=teal!4,colframe=teal!35!black,title={Before you move on}]
What does \mr{नाणेफेक} mean? (\textquotedblleft{}A coin toss\textquotedblright{}.) Say it once more.
\end{tcolorbox}
\section[\texorpdfstring{\mr{पैज} \mr{लावणे} (paij lāvṇe)}{पैज लावणे (paij lāvṇe)}]{\mr{पैज} \mr{लावणे} (paij lāvṇe) — to bet}

\label{lesson:MR-C550-paij-lavne}



\begin{quote}
Before the new one: say the Marathi for good fortune, then the Marathi for a coin toss.
\end{quote}

\subsection*{You'll want to know: \mr{पैज} \mr{लावणे}}
\textbf{\mr{पैज} \mr{लावणे}} — \emph{paij lāvṇe} — \textquotedblleft{}to bet\textquotedblright{}.

\begin{quote}
\mr{कोण} \mr{जिंकेल} \mr{यावर} \mr{दोन} \mr{मित्रांनी} \mr{पैज} \mr{लावली}.
\end{quote}

\begin{quote}
\emph{koṇ jiṅkel yāvar don mitrānnī paij lāvlī.}
\end{quote}

\begin{quote}
Two friends bet on who would win.
\end{quote}

Past, with -\mr{ने} (he …): \mr{त्याने} \mr{पैज} \mr{लावली}, \emph{tyāne paij lāvlī}.

The question takes -\mr{वर}: \mr{यावर} \mr{पैज} \mr{लावणे}.

\begin{cousinweb}
\textbf{\mr{पैज} \mr{लावणे}} (\emph{paij lāvṇe}) is \mr{पैज}, a tadbhava: CDIAL lists Marathi \emph{paij}, \textquotedblleft{}bet, promise\textquotedblright{}, under Sanskrit \emph{pratijñā}, \textquotedblleft{}promise\textquotedblright{} (CDIAL 8563), with \mr{लावणे}.
\end{cousinweb}

\begin{grammarlens}[title={What you've built}]
One more word for saying what happens if.
\end{grammarlens}

\subsection*{Your turn}
\begin{itemize}
\raggedright
  \item \emph{Say it:} \emph{paij lāvṇe}
  \item \emph{Say it:} \emph{paij lāvṇe}, once more
  \item \emph{Say it:} \emph{koṇ jiṅkel yāvar don mitrānnī paij lāvlī.} — Two friends bet on who would win.
  \item \emph{Recall:} say the Marathi for good fortune, then the Marathi for a coin toss, then say \emph{paij lāvṇe} again
\end{itemize}

\begin{tcolorbox}[breakable,colback=teal!4,colframe=teal!35!black,title={Before you move on}]
What does \mr{पैज} \mr{लावणे} mean? (\textquotedblleft{}To bet\textquotedblright{}.) Say it once more.
\end{tcolorbox}
\section[(dialogue)]{Looking back: from Outlines and Drafts to Luck and Chance}

\label{lesson:MR-R550-review-544-550}



\begin{quote}
Say the Marathi for an outline and a draft.
\end{quote}

\subsection*{Your turn}
Say these aloud:

\begin{itemize}
\raggedright
  \item an outline, a draft, to plan ahead, to postpone, to see through
  \item destiny, a golden chance, a convenience, to conjecture, to forecast
  \item a resolution, determination, an intention, to take an oath, to dare
  \item an emergency, a fire extinguisher, an alarm, a security guard, to bring under control
  \item an obstacle, a way out, to bear a load, to sort out, to waver
  \item to compensate, to miss out on, a shield, protection, reliance
  \item a coincidence, gambling, good fortune, a coin toss, to bet
\end{itemize}

\begin{tcolorbox}[breakable,colback=teal!4,colframe=teal!35!black,title={Before you move on}]
An outline? (\textbf{\mr{आराखडा}}, \emph{ārākhḍā}.) A draft? (\textbf{\mr{मसुदा}}, \emph{masudā}.) To plan ahead? (\textbf{\mr{नियोजन} \mr{करणे}}, \emph{niyojan karṇe}.) To postpone? (\textbf{\mr{लांबणीवर} \mr{टाकणे}}, \emph{lāmbṇīvar ṭākṇe}.) To see through? (\textbf{\mr{तडीस} \mr{नेणे}}, \emph{taḍīs neṇe}.) Destiny? (\textbf{\mr{भवितव्य}}, \emph{bhavitavya}.) A golden chance? (\textbf{\mr{सुवर्णसंधी}}, \emph{suvarṇasandhī}.) A convenience? (\textbf{\mr{सोय}}, \emph{soy}.) To conjecture? (\textbf{\mr{अटकळ} \mr{बांधणे}}, \emph{aṭkaḷ bāndhṇe}.) To forecast? (\textbf{\mr{वर्तवणे}}, \emph{vartavṇe}.) A resolution? (\textbf{\mr{संकल्प}}, \emph{saṅkalpa}.) Determination? (\textbf{\mr{निर्धार}}, \emph{nirdhār}.) An intention? (\textbf{\mr{इरादा}}, \emph{irādā}.) To take an oath? (\textbf{\mr{शपथ} \mr{घेणे}}, \emph{śapath gheṇe}.) To dare? (\textbf{\mr{धजणे}}, \emph{dhajṇe}.) An emergency? (\textbf{\mr{आणीबाणी}}, \emph{āṇībāṇī}.) A fire extinguisher? (\textbf{\mr{अग्निशामक}}, \emph{agniśāmak}.) An alarm? (\textbf{\mr{गजर}}, \emph{gajar}.) A security guard? (\textbf{\mr{सुरक्षारक्षक}}, \emph{surakṣārakṣak}.) To bring under control? (\textbf{\mr{आटोक्यात} \mr{आणणे}}, \emph{āṭokyāt āṇṇe}.) An obstacle? (\textbf{\mr{अडथळा}}, \emph{aḍthaḷā}.) A way out? (\textbf{\mr{तोडगा}}, \emph{toḍgā}.) To bear a load? (\textbf{\mr{पेलणे}}, \emph{pelṇe}.) To sort out? (\textbf{\mr{निस्तरणे}}, \emph{nistarṇe}.) To waver? (\textbf{\mr{डळमळणे}}, \emph{ḍaḷmaḷṇe}.) To compensate? (\textbf{\mr{भरपाई} \mr{देणे}}, \emph{bharpāī deṇe}.) To miss out on? (\textbf{\mr{मुकणे}}, \emph{mukṇe}.) A shield? (\textbf{\mr{कवच}}, \emph{kavac}.) Protection? (\textbf{\mr{संरक्षण}}, \emph{saṁrakṣaṇ}.) Reliance? (\textbf{\mr{भरवसा}}, \emph{bharavsā}.) A coincidence? (\textbf{\mr{योगायोग}}, \emph{yogāyog}.) Gambling? (\textbf{\mr{जुगार}}, \emph{jugār}.) Good fortune? (\textbf{\mr{भाग्य}}, \emph{bhāgya}.) A coin toss? (\textbf{\mr{नाणेफेक}}, \emph{nāṇephek}.) To bet? (\textbf{\mr{पैज} \mr{लावणे}}, \emph{paij lāvṇe}.)
\end{tcolorbox}
